\section{What the criterion says about schedules}
\label{sec:solvers}
In additive noise coordinates the probability-flow drift is $f(x,s)=(x-D_s(x))/s$~\cite{karras2022}. Write $\rho=b/a\in(0,1)$ for the denoising ratio. Euler's update is $\rho x+(1-\rho)D_a(x)$, a strongly monotone map with lower slope bound $\rho$ and one denoiser call. The two-evaluation methods in \cref{tab:coeff} have the same residual coefficient $r=\rho$ and differ in how they combine predictions.

For \textbf{sigma-Heun}, the Euler predictor at $b$ is $p_a=\rho x+(1-\rho)D_a(x)$. Substituting it into the trapezoidal update $x+(b-a)[f(x,a)+f(p_a,b)]/2$ produces $w_1$ with the factor $2\rho-1$. The step is universally injective exactly for $\rho\ge1/2$.
For \textbf{explicit sigma-midpoint}, the stage is an Euler step to $(a+b)/2$. Inserting this stage into the midpoint drift gives a negative $w_1$ for every strict denoising ratio.
For \textbf{unthresholded log-SNR midpoint DPM++2S}, the intermediate noise is $\sqrt{ab}$ and the final update uses only the prediction there~\cite{lu2022}. Thus $w_1=0$ and $w_2=1-\rho$. All strict positive denoising ratios satisfy the sign test. Supplement~B derives the three formulas in the same coordinates. When $b=a$, each of these solver steps is the identity.

\begin{table}[t]
\centering\small
\setlength{\tabcolsep}{3.5pt}
\begin{tabular}{@{}lccc@{}}
\toprule
Step & $q=\kappa$ & $w_1$ & $w_2$\\
\midrule
Heun & $\rho$ & $\dfrac{(1-\rho)(2\rho-1)}{2\rho}$ & $\dfrac{1-\rho}{2\rho}$\\[7pt]
Sigma-midpoint & $\dfrac{1+\rho}{2}$ & $-\dfrac{(1-\rho)^2}{1+\rho}$ & $\dfrac{2(1-\rho)}{1+\rho}$\\[7pt]
DPM++2S & $\sqrt{\rho}$ & $0$ & $1-\rho$\\
\bottomrule
\end{tabular}
\caption{Prediction weights in \cref{eq:family}, with $r=\rho=b/a\in(0,1)$. The sign condition applies to the specified ordinary steps with exact posterior means.}
\label{tab:coeff}
\end{table}

\begin{corollary}[Fixed Heun schedules]
\label{cor:schedule}
Fix the ratios $\rho_j\in(0,1]$ of an input-independent denoising schedule. The standard sigma-Heun composition is a global orientation-preserving smooth diffeomorphism for every dimension, common compact prior and positive global scale exactly when every $\rho_j\ge1/2$. For overall $0<b<a$, the minimum length of a schedule with this guarantee is
\begin{equation}
 n_{\min}=\big\lceil\log_2(a/b)\big\rceil.
 \label{eq:nmin}
\end{equation}
\end{corollary}
\begin{proof}
When all ratios meet the threshold, the composition inherits the guarantee. Otherwise use the \emph{first} ratio below $1/2$. With the fixed two-atom prior, reducing the global scale makes that step's source noise small enough to produce a collision. The preceding steps are onto, so two distinct initial states reach the two colliding inputs. Every later deterministic step acts on their equal state and preserves the collision.

The product of $n$ ratios each at least $1/2$ is at least $2^{-n}$. Conversely, $n$ equal ratios $(b/a)^{1/n}$ meet the threshold whenever $n\ge\lceil\log_2(a/b)\rceil$.
\end{proof}

At $2\to1/8$, the grid $2,1,1/2,1/4,1/8$ uses the minimum four Heun substeps and eight denoiser calls. One DPM++2S step has the universal guarantee with two calls. Euler illustrates the separate role of forward fidelity even with a one-call guarantee. The Heun limit is a structural budget for the specified family of schedules. Numerical accuracy and the work of solving for a source must be measured separately.
